%% Generated by Sphinx.
\def\sphinxdocclass{report}
\documentclass[letterpaper,10pt,spanish]{sphinxmanual}
\ifdefined\pdfpxdimen
   \let\sphinxpxdimen\pdfpxdimen\else\newdimen\sphinxpxdimen
\fi \sphinxpxdimen=.75bp\relax
\ifdefined\pdfimageresolution
    \pdfimageresolution= \numexpr \dimexpr1in\relax/\sphinxpxdimen\relax
\fi
\newdimen\sphinxremdimen\sphinxremdimen = 10pt
%% let collapsible pdf bookmarks panel have high depth per default
\PassOptionsToPackage{bookmarksdepth=5}{hyperref}

\PassOptionsToPackage{booktabs}{sphinx}
\PassOptionsToPackage{colorrows}{sphinx}

\PassOptionsToPackage{warn}{textcomp}
\usepackage[utf8]{inputenc}
\ifdefined\DeclareUnicodeCharacter
% support both utf8 and utf8x syntaxes
  \ifdefined\DeclareUnicodeCharacterAsOptional
    \def\sphinxDUC#1{\DeclareUnicodeCharacter{"#1}}
  \else
    \let\sphinxDUC\DeclareUnicodeCharacter
  \fi
  \sphinxDUC{00A0}{\nobreakspace}
  \sphinxDUC{2500}{\sphinxunichar{2500}}
  \sphinxDUC{2502}{\sphinxunichar{2502}}
  \sphinxDUC{2514}{\sphinxunichar{2514}}
  \sphinxDUC{251C}{\sphinxunichar{251C}}
  \sphinxDUC{2572}{\textbackslash}
\fi
\usepackage{cmap}
\usepackage[T1]{fontenc}
\usepackage{amsmath,amssymb,amstext}
\usepackage{babel}



\usepackage{tgtermes}
\usepackage{tgheros}
\renewcommand{\ttdefault}{txtt}



\usepackage[Sonny]{fncychap}
\ChNameVar{\Large\normalfont\sffamily}
\ChTitleVar{\Large\normalfont\sffamily}
\usepackage{sphinx}

\fvset{fontsize=auto}
\usepackage{geometry}


% Include hyperref last.
\usepackage{hyperref}
% Fix anchor placement for figures with captions.
\usepackage{hypcap}% it must be loaded after hyperref.
% Set up styles of URL: it should be placed after hyperref.
\urlstyle{same}

\addto\captionsspanish{\renewcommand{\contentsname}{Contents:}}

\usepackage{sphinxmessages}
\setcounter{tocdepth}{1}



\title{comparador ISS}
\date{10 de abril de 2026}
\release{es}
\author{ismael}
\newcommand{\sphinxlogo}{\vbox{}}
\renewcommand{\releasename}{Versión}
\makeindex
\begin{document}

\ifdefined\shorthandoff
  \ifnum\catcode`\=\string=\active\shorthandoff{=}\fi
  \ifnum\catcode`\"=\active\shorthandoff{"}\fi
\fi

\pagestyle{empty}
\sphinxmaketitle
\pagestyle{plain}
\sphinxtableofcontents
\pagestyle{normal}
\phantomsection\label{\detokenize{index::doc}}


\sphinxAtStartPar
Add your content using \sphinxcode{\sphinxupquote{reStructuredText}} syntax. See the
\sphinxhref{https://www.sphinx-doc.org/en/master/usage/restructuredtext/index.html}{reStructuredText}
documentation for details.


\chapter{Comparador ISS}
\label{\detokenize{index:comparador-iss}}
\sphinxstepscope


\section{Ejercicio 5\_3}
\label{\detokenize{modules:ejercicio-5-3}}\label{\detokenize{modules::doc}}
\sphinxstepscope


\subsection{main module}
\label{\detokenize{main:module-main}}\label{\detokenize{main:main-module}}\label{\detokenize{main::doc}}\index{module@\spxentry{module}!main@\spxentry{main}}\index{main@\spxentry{main}!module@\spxentry{module}}
\sphinxAtStartPar
Actividad 5\_3 \sphinxhyphen{} Comparación de optimización y documentación en Python

\sphinxAtStartPar
Aplicación con interfaz gráfica (Tkinter) que compara el rendimiento
entre una versión optimizada y una no optimizada utilizando datos en
tiempo real de la ISS (International Space Station).

\sphinxAtStartPar
Incluye:
\sphinxhyphen{} Consumo de API
\sphinxhyphen{} Comparación de rendimiento
\sphinxhyphen{} Perfilado con cProfile
\sphinxhyphen{} Interfaz gráfica con Tkinter
\index{App (clase en main)@\spxentry{App}\spxextra{clase en main}}

\begin{fulllineitems}
\phantomsection\label{\detokenize{main:main.App}}
\pysigstartsignatures
\pysiglinewithargsret
{\sphinxbfcode{\sphinxupquote{\DUrole{k}{class}\DUrole{w}{ }}}\sphinxcode{\sphinxupquote{main.}}\sphinxbfcode{\sphinxupquote{App}}}
{\sphinxparam{\DUrole{n}{root}}}
{}
\pysigstopsignatures
\sphinxAtStartPar
Bases: \sphinxcode{\sphinxupquote{object}}

\sphinxAtStartPar
Clase principal de la aplicación Tkinter.

\sphinxAtStartPar
Gestiona la interfaz gráfica, la comparación entre versiones
y la actualización en tiempo real de los datos.
\index{crear\_panel() (método de main.App)@\spxentry{crear\_panel()}\spxextra{método de main.App}}

\begin{fulllineitems}
\phantomsection\label{\detokenize{main:main.App.crear_panel}}
\pysigstartsignatures
\pysiglinewithargsret
{\sphinxbfcode{\sphinxupquote{crear\_panel}}}
{\sphinxparam{\DUrole{n}{frame}}\sphinxparamcomma \sphinxparam{\DUrole{n}{titulo}}\sphinxparamcomma \sphinxparam{\DUrole{n}{funcion}}}
{}
\pysigstopsignatures
\sphinxAtStartPar
Crea un panel de comparación en la interfaz.
\begin{quote}\begin{description}
\sphinxlineitem{Parámetros}\begin{itemize}
\item {} 
\sphinxAtStartPar
\sphinxstyleliteralstrong{\sphinxupquote{frame}} (\sphinxstyleliteralemphasis{\sphinxupquote{tk.Frame}}) \textendash{} Contenedor del panel.

\item {} 
\sphinxAtStartPar
\sphinxstyleliteralstrong{\sphinxupquote{titulo}} (\sphinxstyleliteralemphasis{\sphinxupquote{str}}) \textendash{} Título del panel.

\item {} 
\sphinxAtStartPar
\sphinxstyleliteralstrong{\sphinxupquote{funcion}} (\sphinxstyleliteralemphasis{\sphinxupquote{callable}}) \textendash{} Función de obtención de datos.

\end{itemize}

\end{description}\end{quote}

\end{fulllineitems}

\index{dibujar\_estrellas() (método de main.App)@\spxentry{dibujar\_estrellas()}\spxextra{método de main.App}}

\begin{fulllineitems}
\phantomsection\label{\detokenize{main:main.App.dibujar_estrellas}}
\pysigstartsignatures
\pysiglinewithargsret
{\sphinxbfcode{\sphinxupquote{dibujar\_estrellas}}}
{}
{}
\pysigstopsignatures
\sphinxAtStartPar
Dibuja un fondo de estrellas aleatorias para estética espacial.

\end{fulllineitems}

\index{mostrar\_help() (método de main.App)@\spxentry{mostrar\_help()}\spxextra{método de main.App}}

\begin{fulllineitems}
\phantomsection\label{\detokenize{main:main.App.mostrar_help}}
\pysigstartsignatures
\pysiglinewithargsret
{\sphinxbfcode{\sphinxupquote{mostrar\_help}}}
{}
{}
\pysigstopsignatures
\sphinxAtStartPar
Muestra ayuda con la documentación de las funciones principales.

\end{fulllineitems}


\end{fulllineitems}

\index{obtener\_datos\_no\_opt() (en el módulo main)@\spxentry{obtener\_datos\_no\_opt()}\spxextra{en el módulo main}}

\begin{fulllineitems}
\phantomsection\label{\detokenize{main:main.obtener_datos_no_opt}}
\pysigstartsignatures
\pysiglinewithargsret
{\sphinxcode{\sphinxupquote{main.}}\sphinxbfcode{\sphinxupquote{obtener\_datos\_no\_opt}}}
{}
{}
\pysigstopsignatures
\sphinxAtStartPar
Obtiene datos de la ISS de forma NO optimizada.

\sphinxAtStartPar
Simula un código ineficiente con bucles innecesarios y cálculos redundantes.
\begin{quote}\begin{description}
\sphinxlineitem{Devuelve}
\sphinxAtStartPar
Diccionario con latitud y longitud promedio.

\sphinxlineitem{Tipo del valor devuelto}
\sphinxAtStartPar
dict

\end{description}\end{quote}

\end{fulllineitems}

\index{obtener\_datos\_opt() (en el módulo main)@\spxentry{obtener\_datos\_opt()}\spxextra{en el módulo main}}

\begin{fulllineitems}
\phantomsection\label{\detokenize{main:main.obtener_datos_opt}}
\pysigstartsignatures
\pysiglinewithargsret
{\sphinxcode{\sphinxupquote{main.}}\sphinxbfcode{\sphinxupquote{obtener\_datos\_opt}}}
{}
{}
\pysigstopsignatures
\sphinxAtStartPar
Obtiene datos de la ISS de forma optimizada.

\sphinxAtStartPar
Accede directamente a la API sin estructuras innecesarias.
\begin{quote}\begin{description}
\sphinxlineitem{Devuelve}
\sphinxAtStartPar
Diccionario con latitud y longitud actuales.

\sphinxlineitem{Tipo del valor devuelto}
\sphinxAtStartPar
dict

\end{description}\end{quote}

\end{fulllineitems}

\index{perfilar() (en el módulo main)@\spxentry{perfilar()}\spxextra{en el módulo main}}

\begin{fulllineitems}
\phantomsection\label{\detokenize{main:main.perfilar}}
\pysigstartsignatures
\pysiglinewithargsret
{\sphinxcode{\sphinxupquote{main.}}\sphinxbfcode{\sphinxupquote{perfilar}}}
{\sphinxparam{\DUrole{n}{func}}}
{}
\pysigstopsignatures
\sphinxAtStartPar
Ejecuta una función y mide su rendimiento con cProfile.
\begin{quote}\begin{description}
\sphinxlineitem{Parámetros}
\sphinxAtStartPar
\sphinxstyleliteralstrong{\sphinxupquote{func}} (\sphinxstyleliteralemphasis{\sphinxupquote{function}}) \textendash{} Función a analizar.

\sphinxlineitem{Devuelve}
\sphinxAtStartPar
\begin{itemize}
\item {} 
\sphinxAtStartPar
dict: Resultado de la función

\item {} 
\sphinxAtStartPar
float: Tiempo de ejecución

\item {} 
\sphinxAtStartPar
str: Estadísticas de profiling (cProfile)

\end{itemize}


\sphinxlineitem{Tipo del valor devuelto}
\sphinxAtStartPar
tuple

\end{description}\end{quote}

\end{fulllineitems}



\renewcommand{\indexname}{Índice de Módulos Python}
\begin{sphinxtheindex}
\let\bigletter\sphinxstyleindexlettergroup
\bigletter{m}
\item\relax\sphinxstyleindexentry{main}\sphinxstyleindexpageref{main:\detokenize{module-main}}
\end{sphinxtheindex}

\renewcommand{\indexname}{Índice}
\printindex
\end{document}